% 1. Introduction
\section{Introduction}

Whole-brain connectomes now make it possible to ask, at synaptic
resolution, which individual neurons matter most for the directed flow of
information through an entire central nervous system. For the adult female
\emph{Drosophila melanogaster}, the FlyWire reconstruction of the Full
Adult Fly Brain (FAFB) volume provides a dense, verified wiring diagram of
139,255 neurons and roughly 3.7 million thresholded synaptic connections
\cite{dorkenwald2024}, with cell typing and neurotransmitter predictions
\cite{schlegel2024} and a complete visual-system parts list
\cite{matsliah2024}.

Directed network structure matters because the brain is not an undirected
web of interactions: influence propagates along asymmetric anatomical
paths, and the global ease of communication --- formalized as network
efficiency, the mean inverse shortest-path length over all directed neuron
pairs --- depends on how those paths are organized. Prior network analyses
of the fly connectome established rich-club organization, motif structure,
and the important observation that high-total-degree neurons are mostly
GABAergic \cite{lin2024}; earlier work examined information flow at coarser
resolution \cite{shih2015}, and hierarchical modularity has been
characterized \cite{kunin2023}.

A neuron's importance as a \emph{structural control point}, however, cannot
be read from its degree alone. Bridging positions can make sparse neurons
disproportionately influential, whereas removing one of several parallel
hubs may barely affect global communication. Removal-based perturbation ---
delete a node, recompute global efficiency, and record the fractional loss
--- captures this directly. Such analyses exist in aggregate,
degree-ordered form for the fly brain \cite{lin2024} and in other systems
\cite{uzel2022}, but a per-neuron, neurotransmitter-resolved,
degree-controlled ranking with null-model validation has not, to our
knowledge, been reported for the adult connectome.

Neurotransmitter identity is a natural candidate organizing variable: if
inhibition were preferentially positioned at high-leverage network
locations, this would constrain how the brain balances excitation and
control. The association between GABAergic identity and high degree is
already documented \cite{lin2024}; the open question is whether any
\emph{excess} control impact remains after degree is accounted for.

We therefore pre-registered a perturbation pipeline and froze the
hypothesis before unblinding: \emph{after accounting for degree, GABAergic
identity is associated with disproportionately high per-neuron structural
control impact.} We tested this with (i) neuron-level removal
perturbation on the FAFB v783 connectome using a validated Control Impact
Score (CIS); (ii) degree-matched GABA--ACh comparison; (iii) 100
degree-preserving null networks; and (iv) a regression of control impact
on degree and neurotransmitter identity. As a pre-planned secondary
analysis, we mapped where the strongest structural chokepoints reside and
tested their anatomical enrichment against degree-matched controls.

Throughout, we distinguish network topology from functional neural
dynamics: every quantity here is structural, computed from the wiring
diagram alone. ``Structural control impact'' denotes sensitivity of
global efficiency to computational node removal; it does not measure
firing, causality, or behavioral necessity.
